\SPBGContentSlide
  {Реальные данные: парные проверки}
  {%
    \centering
    \TinyBody
    \renewcommand{\arraystretch}{1.12}
    \setlength{\tabcolsep}{2.6pt}
    \begin{tabular}{|l|r|r|c|r|r|c|}
      \hline
      \textbf{Метод} & \textbf{F1} & \textbf{$\Delta F_1$} & \textbf{$p_H$}
      & \textbf{$t$, c} & \textbf{$\Delta t$} & \textbf{$p_H$} \\
      \hline
      MDM-G & 0.823 & +0.0003 & 0.379 & 3.1236 & +1.6113 & $8.87\cdot10^{-16}$ \\
      \hline
      MDM-A & 0.822 & -0.0005 & 0.613 & 6.3971 & +4.8848 & $8.87\cdot10^{-16}$ \\
      \hline
      SMO & 0.815 & -0.007 & 0.613 & 0.2357 & -1.2766 & $8.87\cdot10^{-16}$ \\
      \hline
      SVC & 0.844 & +0.021 & 0.017 & 0.1988 & -1.3135 & $8.87\cdot10^{-16}$ \\
      \hline
      LinearSVC & 0.846 & +0.024 & 0.053 & 0.0027 & -1.5096 & $8.87\cdot10^{-16}$ \\
      \hline
    \end{tabular}

    \vspace{.18cm}
    \begin{minipage}{.90\linewidth}
      \SmallBody
      \raggedright
      На R1--R15 различия MDM, MDM-G и MDM-A по качеству не подтверждаются.
      SVC значимо превосходит MDM по F1 и ROC-AUC; у LinearSVC F1 выше, но
      после поправки Холма $p_H=0.053$. По времени все отличия от MDM значимы.
    \end{minipage}
  }
